\section{Introduction}

Text embedding models have become a basic infrastructure layer for semantic search, clustering, reranking, and retrieval-augmented generation. Their usefulness comes not only from individual vectors, but from the geometric structure those vectors induce: neighborhoods, semantic regions, and smooth paths between related examples. Strong recent embedding models, including LLM-based encoders, often learn rich semantic geometry but are costly to run at scale. Encoding large private collections, refreshing indexes as documents change, or serving retrieval on constrained hardware therefore motivates compact student encoders that inherit the geometry of a stronger teacher.

Knowledge distillation offers a practical route to this compression~\citep{hinton2015distilling}. In language representation learning, distillation has been applied through output matching, hidden-state alignment, self-attention transfer, and compact pretrained architectures~\citep{sanh2019distilbert,jiao2019tinybert,wang2020minilm,xu2024survey}. For embedding models, common objectives include pointwise alignment between student and teacher embeddings, contrastive learning with teacher-derived positives, teacher-anchored layer alignment, and relational losses that match pairwise similarities within a mini-batch. These objectives provide valuable supervision at different granularities: individual examples, internal layers, sampled pairs, or local batch structure.

The geometric view suggests a complementary form of supervision. A teacher embedding space can be treated as a sampled semantic manifold: nearby examples define local neighborhoods and multi-hop paths reveal semantic continuity. Where pointwise losses encourage coordinate agreement, layer alignment improves internal transfer under capacity mismatch, and mini-batch relational losses preserve sampled similarity patterns, this view raises a different question: can a compact student inherit the teacher's semantic manifold by matching heat-style diffusion on a teacher-induced neighborhood graph? If the teacher embeddings define a discrete manifold over training examples, random walks on it reveal how semantic mass flows from each anchor to nearby examples, broader regions, and higher-order communities. Matching that diffusion process gives the student geometry-aware supervision tied neither to teacher internals nor to a particular student architecture.

We propose \textbf{HeatGeo}, a heat-diffusion manifold distillation method for compact text embedding models. HeatGeo first caches teacher embeddings and constructs a mutual $k$-nearest-neighbor graph using teacher cosine similarity. This graph acts as a discrete approximation to the teacher's semantic manifold. From it, HeatGeo defines a transition matrix and precomputes diffusion targets across a range of diffusion times, from the ambient teacher metric itself to one-step, two-step, and longer-range semantic random walks. During training, each anchor example is paired with a candidate set containing high-diffusion neighbors, teacher hard negatives, and random negatives, and candidate sets are pooled across the mini-batch so that anchors are compared on a common set of columns. The student then learns neighborhood distributions over those columns and minimizes a weighted KL divergence to the teacher targets. A second term extends supervision beyond the anchors: random walks on the teacher graph select non-anchor nodes that the batch has already encoded, and the student is trained to reproduce the teacher's one-step kernel at those nodes as well, which both densifies supervision and pins down the similarity offsets that anchor-only matching leaves free.

A single design decision makes this family more than target smoothing: the student matches each scale at its own resolution, which keeps the scales distinct constraints on the student instead of collapsing them into one smoothed target (Proposition~\ref{prop:collapse}). HeatGeo needs no pointwise regression, contrastive pair construction, layer alignment, or graph-coordinate objective, and adds no trainable parameters beyond the student encoder.

Our contributions are:
\begin{itemize}
    \item We formulate text embedding distillation as manifold geometry transfer, using the teacher embedding space as a discrete semantic manifold.
    \item We develop a theory that derives each component of the objective from the failure mode of batch-relational supervision: why supervision must be anchored to local neighborhoods, why matching them locally controls global geometry, why small sampled candidate sets suffice, and why one calibration term completes the loss (Section~\ref{sec:theory}).
    \item We introduce HeatGeo, which distills a diffusion-time family of heat kernels from a teacher-induced mutual $k$NN graph into a compact student encoder, with resolution-dependent student temperatures that keep the scales informative and no requirement for teacher hidden states or task labels.
\end{itemize}
